% ============================================================================
% 05-resultados.tex — Resultados (módulo 6/8)
% ============================================================================
\chapter{Resultados}

\section{Ranking dos modelos free (R499)}

A Tabela \ref{tab:rank} consolida o \textit{benchmark} de modelos free reais
na condição decisiva C2 (\textit{scaffold} MASWOS), com latências medidas e
acurácia (correção do par $(7,3)$), conformidade de formato e \textit{score}
composto. Na condição C0, ambos os modelos testados retornaram vazio; na C1,
o \texttt{nemotron} excedeu o \textit{timeout} de 110~s. O \textit{scaffold}
MASWOS foi decisivo: 0\% $\rightarrow$ 100\% (primeiro \textit{trial}, modelos
que resolveram).

\begin{table}[htbp]
\centering
\caption{Ranking dos modelos free na condição C2 (MASWOS), R499.}
\label{tab:rank}
\begin{tabular}{lccccl}
\toprule
Modelo & Latência (s) & Acurácia & Conformidade & Score & Status \\
\midrule
\texttt{mimo-v2.5-free}            & 23,8  & 1,0 & 1,0 & \textbf{0,987} & corrigiu, typo unicode cosmético \\
\texttt{big-pickle}                & 43,0  & 1,0 & 1,0 & 0,923 & corrigiu e conforme, sem typo \\
\texttt{nemotron-3-ultra-free}     & 107,4 & 1,0 & 1,0 & 0,709 & corrigiu; 4--5$\times$ mais lento \\
\texttt{muse-spark-1.3-free}       & 18,1  & 0,0 & 0,0 & 0,324 & ``fake reasoning'': anuncia verificação, sem resposta \\
\texttt{muse-spark-1.2-free}       & 24,4  & 0,0 & 0,0 & 0,312 & idem \\
\bottomrule
\end{tabular}
\end{table}

\textbf{Achados de infraestrutura (R499)}: \texttt{deepseek-v4-flash} não
disponível (saldo insuficiente); \texttt{gemini-2.5-flash} retornou erro de
servidor; \texttt{gpt-4o-mini} e \texttt{claude-haiku-4} sem resposta. O
tier \texttt{fast} do catálogo não é sinônimo de acessível — limitação de
catálogo mapeada.

\section{Replicação pós-R500 e não-determinismo}

A Tabela \ref{tab:repl} resume a replicação conduzida após a implementação do
roteamento R500.

\begin{table}[htbp]
\centering
\caption{Replicação pós-R500 da condição C2 via roteamento automático.}
\label{tab:repl}
\begin{tabular}{lcccl}
\toprule
Modelo & Rota & n & Acurácia replicada & Observação \\
\midrule
\texttt{mimo-v2.5-free}   & \texttt{route\_free("academic")} & 3 & 2/3 (0,67) & 1 timeout de infra (110~s) \\
\texttt{big-pickle}       & \texttt{route\_free("writing")}  & 2 & 0/2 (0,00) & output vazio; hipótese: estado do host sem reinício \\
\bottomrule
\end{tabular}
\end{table}

Os resultados confirmam a estocasticidade do \textit{decoding}: o mesmo prompt
C2 acertou (26,6~s; 20,9~s) e errou por \textit{timeout} (110~s) no
\texttt{mimo}; o \texttt{big-pickle} retornou vazio nas duas réplicas
(20,9~s; 41,7~s), contra acerto C2 de 43,0~s no R499. A hipótese de estado
acumulado do \textit{host} opencode (sessão de 1h15m sem reinício) foi
registrada, mas o reinício não pôde ser executado sem encerrar a sessão
ativa. \textbf{Decisão metodológica}: \texttt{big-pickle} é reportado como
\textit{score} pré-R500 não confirmado na replicação; o \texttt{mimo} é
reportado com média n=3 (0,67), não 1,0. Esta é a principal contribuição
anti-\textit{overclaim} do ciclo R500: \textbf{1 \textit{trial} por condição
não estima a taxa real de LLMs}.

\section{Roteamento automático (R500)}

O fluxograma da Figura \ref{fig:flux-roteamento} documenta a decisão
automatizada: \texttt{route(``academic'', prefer\_free=True)} retorna
\texttt{opencode/mimo-v2.5-free} com \texttt{reason} indicando curadoria R499
(score 0,987) e alternativas ordenadas (\texttt{big-pickle}, \texttt{nemotron},
\texttt{muse-1.3}); \texttt{list\_all\_models()} inclui o catálogo free.
Regressão: perfis existentes inalterados sem \texttt{prefer\_free} (46 testes
verdes, 1,97~s).

% ============================================================================
% fig-flux-roteamento.tex — Fluxograma do roteamento R500 (Figura 2)
% ============================================================================
\begin{figure}[htbp]
\centering
\begin{tikzpicture}[
  node distance=0.85cm and 1.3cm,
  caixa/.style={rectangle, draw=cororo, thick, rounded corners=2pt,
                fill=cororo!8, minimum width=3.1cm, minimum height=0.9cm,
                align=center, font=\footnotesize},
  decisao/.style={diamond, draw=coramarelo, thick, fill=coramarelo!10,
                  aspect=2, align=center, font=\footnotesize},
  seta/.style={-{Stealth[length=3mm]}, thick, color=corcinza}
]
  \node[caixa] (tarefa) {Tarefa acadêmica\\ \texttt{route(task\_type, prefer\_free=True)}};
  \node[decisao, right=of tarefa] (dec) {task\_type\\ em \{academic, math,\\ reasoning, writing\}?};
  \node[caixa, below=of dec] (cur) {Curadoria R499\\ \texttt{free\_model\_catalog.py}};
  \node[caixa, below=of cur] (mimo) {melhor score\\ \texttt{mimo-v2.5-free} (0,987)};
  \node[caixa, right=of cur] (alt) {fallbacks: big-pickle (0,923)\\ nemotron (0,709)\\ muse-1.3 (0,324)};
  \node[caixa, below=of mimo] (rota) {\texttt{RouteResult} com\\ \texttt{reason}=``Curadoria R499''};

  \draw[seta] (tarefa) -- (dec);
  \draw[seta] (dec) -- node[right, font=\tiny]{sim} (cur);
  \draw[seta] (cur) -- (mimo);
  \draw[seta] (mimo) -- (rota);
  \draw[seta, dashed, color=corvermelho] (dec) -- ++(2.7,0) node[right, font=\tiny, color=corvermelho]{não: perfis clássicos inalterados};
  \draw[seta, dashed] (cur.east) -- (alt.west);
\end{tikzpicture}
\caption{Fluxograma do roteamento automático R500: a curadoria empírica do
R499 (score composto) decide o modelo free por tipo de tarefa, com cadeia de
\textit{fallback} explícita e \texttt{reason} auditável.}
\label{fig:flux-roteamento}
\end{figure}

\section{Gate SDD/TDD}

O ciclo R500 adicionou 7 testes de roteamento/free-catálogo (verdes); a
regressão \texttt{test\_opencode\_go\_zen.py} segue verde. A suíte completa
pré-R500 registrou 4066 passed, 58 skipped, 0 failed (681~s). O
\texttt{doctor} do ecossistema valida specs, registro de evolução e memória
metacognitiva.